\section{哈密顿纤维化、联络与和乐}
\label{sec:hamiltonian-fibrations-connections-holonomies}

设 $S$ 是亏格为 $g$、具有 $h\geq1$ 个边界分支的紧致、有向、连通曲面：
\[
 \partial S=\bigcup_{i=1}^h\partial_iS,
\]
其中 $\partial_iS$ 是 $S$ 的第 $i$ 个边界分支。$S$ 上的定向在每个
$\partial_iS$ 上诱导边界定向。

考虑带联络 $\Gamma$ 的哈密顿纤维化 $E\to S$，其纤维以
$(M,\omega)$ 为模型。设 $\omega_S$ 是 $S$ 上满足
$\int_S\omega_S=1$ 的面积形式，并设 $\Omega$ 是 $E$ 上与
$\Gamma$ 相伴的耦合形式。

\begin{proposition}\label{prop:weak-coupling-nondegenerate}
存在正常数 $\lambda_0>0$，使得对所有 $\lambda>\lambda_0$，$2$-形式
$\Omega+\lambda\pi^*\omega_S$ 在 $E$ 上非退化。
\end{proposition}

\begin{proof}
我们要证明，存在 $\lambda_0>0$，使得对任意 $\lambda>\lambda_0$，映射
\[
 \xi\longmapsto \xi\mathbin\lrcorner
 (\Omega+\lambda\pi^*\omega_S)
\]
是从 $TE$ 到 $T^*E$ 的同构；也就是说，只要 $\xi$ 满足
\[
 (\Omega+\lambda\pi^*\omega_S)(\xi,\eta)=0
\]
对所有 $\eta\in T_eE$ 成立，就有 $\xi=0$。

设 $\xi$ 满足
$(\Omega+\lambda\pi^*\omega_S)(\xi,\eta)=0$ 对所有
$\eta\in T_eE$ 成立。把
$\xi=\xi^v+\xi^h$、$\eta=\eta^v+\eta^h$ 分解为竖直分量与水平分量，得到
\[
\begin{aligned}
 0&=(\Omega+\lambda\pi^*\omega_S)(\xi,\eta)\\
  &=(\Omega+\lambda\pi^*\omega_S)
    (\xi^v+\xi^h,\eta^v+\eta^h)\\
  &=\Omega(\xi^v,\eta^v)+\Omega(\xi^h,\eta^h)
    +\lambda\pi^*\omega_S(\xi^h,\eta^h),
\end{aligned}
\]
因而
\[
 \Omega(\xi^v,\eta^v)+\Omega(\xi^h,\eta^h)
 =-\lambda\pi^*\omega_S(\xi^h,\eta^h)
\]
对所有 $\eta\in T_eE$ 成立。这里使用了与 $\Omega$ 相容的联络
$\Gamma$ 的定义；该定义蕴含
\[
 \Omega(T^vTM,T^hTM)=0.
\]
分别考虑满足 $\eta^h=0$ 和 $\eta^v=0$ 的 $\eta$，得到
\[
 \begin{cases}
  \Omega(\xi^v,\eta^v)=0,\\
  \Omega(\xi^h,\eta^h)=-\lambda\pi^*\omega_S(\xi^h,\eta^h)
 \end{cases}
\]
对所有 $\eta$ 成立。由 $\Omega$ 的逐纤维非退化性，从第一个等式推出
$\xi^v=0$。另一方面，第二个等式等价于
\[
 \Omega(\xi^h,\eta^h)
 =-\lambda\omega_S(d\pi\xi^h,d\pi\eta^h).
\]
若 $\xi^h=0$，则 $\xi=0$，证明即告完成。否则，对任意满足
$\omega_S(d\pi\xi^h,d\pi\eta^h)\ne0$ 的 $\eta$，上式蕴含
\[
 \frac{|\Omega(\xi^h,\eta^h)|}
 {|\omega_S(d\pi\xi^h,d\pi\eta^h)|}=\lambda.
\]
这样的 $\eta$ 存在，因为 $\omega_S$ 非退化，并且
$d\pi:\Lambda^2(HTE)\to\Lambda^2(TS)$ 是同构。

再者，因为 $S$ 紧，商的上确界
\[
 \sup_{\eta;\,\omega_S(d\pi\xi^h,d\pi\eta^h)\ne0}
 \frac{|\Omega(\xi^h,\eta^h)|}
 {|\omega_S(d\pi\xi^h,d\pi\eta^h)|}
 =
 \sup_{\substack{\xi,\eta;\,
   \omega_S(d\pi\xi^h,d\pi\eta^h)\ne0\\
   |\xi^h|=|\eta^h|=1}}
 \frac{|\Omega(\xi^h,\eta^h)|}
 {|\omega_S(d\pi\xi^h,d\pi\eta^h)|}
\]
非零且有限。把这个值记为 $\lambda_0$。于是，对任意
$\lambda>\lambda_0$ 都会得到矛盾。因此我们已经证明，只要
$\lambda>\lambda_0$，映射
$\xi\mapsto\xi\mathbin\lrcorner(\Omega+\lambda\pi^*\omega_S)$
就是同构。证明完成。
\end{proof}

注意，对所有 $\lambda>\lambda_0$，辛形式
$\Omega_\Gamma+\lambda\pi^*\omega_S$ 彼此形变等价，因而这族形式确定了
$E$ 上由下式给出的辛形式的一个自然形变类：
\[
 \omega_{\Gamma,\lambda}:=\Omega_\Gamma+\lambda\pi^*\omega_S.
\]

\subsection{哈密顿、联络与和乐}
\label{subsec:hamiltonian-connection-holonomy}

下面把和乐效应纳入对有边曲面上哈密顿纤维化的讨论。我们严格遵循
(En00) 的论述，并把讨论限制在 $G=\Ham(M,\omega)$ 的情形。恩托夫
(En00) 是在任意李群 $G$ 的 $G$-丛这一一般背景下，而不只是在哈密顿微分同胚群的背景下，给出这番讨论的。以 $h$ 表示 $\partial S$ 的连通分支数。

注意，在一个平凡化中，哈密顿联络 $\Gamma$ 在每一点
$z\in\partial_iS$ 处都确定了一个位于 $\Ham(M,\omega)$ 中的和乐。换用另一个平凡化，将在 $\Ham(M,\omega)$ 中产生一个共轭元。因此，给定点
$z\in\partial_iS$ 处的和乐在 $\Symp(M,\omega)$ 中确定了
$\Ham(M,\omega)$ 的唯一共轭类。

现在令 $C=(C_1,\ldots,C_h)$ 是 $\Ham(M,\omega)$ 中的一些共轭类，并以
$\mathcal L(C)$ 表示 $E\to S$ 上满足如下条件的哈密顿联络 $\Gamma$ 的集合：
$\Gamma$ 在 $\partial S$ 的某个邻域上平坦，并且它沿 $\partial_iS$ 的和乐属于
$C_i$。由于联络在 $\partial S$ 的一个邻域上平坦，纤维
$E_{b_0}$（$b_0\in\partial_iS$）的平行移动在 $\partial_iS$ 的一个邻域上把
$E$ 平凡化，使得
\[
 E|_U\cong U\times E_{b_0},
 \qquad
 \Gamma_e:T_e(E|_U)=T_{\pi(e)}U\oplus E_{b_0}.
\]
令 $\Omega=\Omega_\Gamma$ 为 $\Gamma$ 的耦合形式。它确定了 $E$ 上一个自然的辛形式形变类，其中的形式为
\[
 \omega_{\Gamma,\lambda}:=\Omega_\Gamma+\lambda\pi^*\omega_S,
\]
这里 $\omega_S$ 是面积形式，而 $\lambda>0$ 是充分大的常数。我们总把
$\omega_S$ 归一化为 $\int_S\omega_S=1$。

\begin{definition}\label{def:hamiltonian-fibration-connection-holonomy}
若三元组 $(E,\Gamma,C)$ 是带联络的哈密顿纤维化 $(E,\Gamma)$，并且它沿
$\partial_iS$ 的和乐 $\phi_i$ 位于 $C_i$ 所包含的一个共轭类中，则称
$(E,\Gamma,C)$ 为带联络与和乐的哈密顿纤维化。
\end{definition}

当 $(g,h)=(0,1)$ 时，$S=D^2$，并且任意哈密顿纤维化
$E\to D^2$ 都是平凡的。设 $(E,\Gamma,C)$ 是带联络并且沿
$\partial D^2$ 具有和乐的哈密顿纤维化。选取平凡化
$\Phi:E\to D^2\times(M,\omega)$，并把 $C$ 表示为
$C=\{[\phi]\}$，其中 $\phi\in\Ham(M,\omega)$。于是，可用一个
$1$-周期哈密顿量 $H:S^1\times M\to\mathbb R$ 表示联络
$\Gamma|_{\partial D^2}$，使得 $[\phi_H^1]\in C$，并且
\[
 (T_e^hE)=
 \left\{
  \left.\left(a\frac{\partial}{\partial t},\xi\right)
  \in T_z\partial D^2\times T_xM\ \right|\
  e=(z,x)\in D^2\times M,\ \xi=aX_H(x)
 \right\}.
\]
等价地，$\partial D^2$ 上的联络 $1$-形式 $A$ 可写为
\[
 A=X_H\,dt.
\]
可以把 $H$ 延拓为一个族 $D^2\times M\to\mathbb R$，仍记作
$H=H(z,x)$。我们要求这个延拓满足
\[
 H(r,t,x)\equiv H(t,x)
\]
对所有 $r\in(1-\delta,1]$ 成立，其中 $\delta>0$ 是某个小数；由于
$\Gamma$ 在 $\partial D^2$ 的一个邻域上平坦，这一点可以做到。更一般地，对任意 $(g,h)$，都可以在 $S$ 上选取一个哈密顿联络
$\Gamma\in\mathcal L(C)$，使得相伴和乐属于给定的共轭类
$C=(C_1,\ldots,C_h)$。

下面的定义由 (En00) 引入。

\begin{definition}\label{def:size-hamiltonian-fibration}
\stfhIndex{I043}{00-symbol-043}{$\operatorname{size}(E)$}
考虑哈密顿纤维化 $E\to S$。对 $E$ 的一个给定哈密顿联络
$\Gamma$，把 $\operatorname{size}(\Gamma)$ 定义为命题
4.3.1 中出现的那些 $\lambda_0$ 的 $1/\lambda_0$ 的上确界，其中形式
$\Omega_\Gamma+\lambda\pi^*\omega_\Sigma$ 非退化。定义
\[
 \operatorname{size}(E):=\sup_\Gamma\operatorname{size}(\Gamma),
\]
其中上确界取遍 $E$ 的所有哈密顿联络 $\Gamma$。
\end{definition}

命题 4.3.1 蕴含 $\operatorname{size}(\Gamma)$ 总是有限的。下一节中，我们将用哈密顿纤维化 $E\to\Sigma$ 的 $K$-面积来估计这个大小。

与此同时，我们希望把上面的讨论与亏格为 $g$、带有 $k$ 个标记点
$\{z_1,\ldots,z_k\}$ 的黎曼曲面 $(\Sigma,j)$ 上的哈密顿纤维化联系起来。复结构 $j$ 在 $\Sigma$ 上诱导复定向。记
\[
 \dot\Sigma=\Sigma\setminus\{z_1,\ldots,z_k\},
\]
并在每个穿孔 $z_i$ 的穿孔邻域 $D_i\setminus\{z_i\}$ 上固定解析坐标，把它等同于带有标准复结构的圆柱
$[0,\infty)\times S^1$ 或 $(-\infty,0]\times S^1$，使得每个
$\partial D_i$ 都映到相应半圆柱的边界圆周。第一类穿孔称为\emph{入穿孔}，第二类称为\emph{出穿孔}。一般地，带有解析坐标的穿孔称为\emph{定向穿孔}。\stfhIndex{I101}{ding^xiang^chuan^kong}{定向穿孔（directed puncture）}

对于正穿孔，把该等同记为
\[
 \varphi_i^+:D_i\setminus\{z_i\}\longrightarrow(-\infty,0]\times S^1,
\]
而对于负穿孔，记为
\[
 \varphi_i^-:D_i\setminus\{z_i\}\longrightarrow[0,\infty)\times S^1.
\]
若把 $[0,\infty)\times S^1$ 或 $(-\infty,0]\times S^1$ 与
$D^2\setminus\{0\}$ 共形等同，则假设这一等同延拓到 $D_i$。我们称这样的等同为\emph{有理坐标}。以 $(\tau,t)$ 表示圆柱上的标准圆柱坐标。于是，有理坐标映射 $\varphi^\pm$ 在 $\tau$ 的平移意义下唯一确定。

现在定义 $\dot\Sigma$ 在 $\vec z=\{z_1,\ldots,z_k\}$ 处的\emph{实爆破}：它是一个紧致有边曲面，其边界由 $k$ 个圆周分支组成；每个圆周以
\stfhIndex{I235}{shi^bao^po}{实爆破（real blow-up）}
\[
 \mathbb{RP}^1\cong\mathbb P_{\mathbb R}(T_{z_i}\Sigma)\cong S^1
\]
代替穿孔 $z_i$。把相应的拓扑空间记为
\[
 Bl_{\vec z}^{\mathbb R}(\Sigma)=S.
\]
通过等同
\[
 \operatorname{Int}(Bl_{\vec z}^{\mathbb R}(\Sigma))
 =\Sigma\setminus\{z_1,\ldots,z_k\},
\]
在 $\operatorname{Int}(Bl_{\vec z}^{\mathbb R}(\Sigma))$ 上赋予自然的共形结构。当还给定 $D_i\setminus\{z_i\}$ 上的解析坐标，例如
$\varphi:D_i\setminus\{z_i\}\to[0,\infty)\times S^1$ 时，就有自然的全纯嵌入
\[
 [0,\infty)\times S^1
 \longrightarrow\operatorname{Int}(Bl_{\vec z}^{\mathbb R}(\Sigma)),
\]
它把 $(0,t)\in\{0\}\times S^1$ 映到 $\partial D_i$。这个全纯嵌入在定义域的相位旋转意义下唯一确定。

当我们指定 $Bl_{\vec z}^{\mathbb R}(\Sigma)$ 各边界分支上的和乐时，也把它们称为 $\dot\Sigma$ 的各穿孔周围的和乐。
\stfhIndex{I168}{chuan^kong^zhou^wei^de^he^le}{穿孔周围的和乐（holonomy around the puncture）}

\begin{lemma}\label{lem:horizontal-section-coupling-area}
考虑哈密顿微分同胚 $\phi$ 以及满足 $H\mapsto\phi$ 的哈密顿量
$H$。令 $M_\phi\to S^1$ 为 $\phi$ 的映射圆，并考虑 $M_\phi$ 上以
$2$-形式
\[
 \omega+d(H\,dt)
\]
为耦合形式的哈密顿联络。考虑任意带联络的哈密顿纤维化
$E\to D^2$，其耦合形式为 $\Omega$，并满足 $\Phi_*\Omega$ 限制为上面
$\partial_iS$ 上的联络。考虑任意沿 $\partial D^2$ 水平的截面
$v:D^2\to E$。那么，对任意给定的平凡化
$\Phi:E\to D^2\times(M,\omega)$，都可写成
\[
 \Phi\circ v(z)=(z,w_\Phi(z)),
\]
其中 $w=w_\Phi:D^2\to M$ 是光滑映射，且对某个 $x\in M$，有
\[
 w|_{\partial D^2}(t)=\phi_H^t(x):=z_H^x(t).
\]
此外，
\begin{equation}
 -\int_{D^2}v^*\Omega
 =-\int_{D^2}w^*\omega
  -\int_{\partial D^2}H(t,z_H^x(t))\,dt.
 \label{eq:horizontal-section-coupling-area}
\end{equation}
\end{lemma}

\begin{proof}
设 $(E,\Omega)$ 是这样的带联络纤维化，其耦合形式为 $\Omega$。在平凡化
$\Phi:E\to D^2\times(M,\omega)$ 中，可写
$\Phi_*\Omega=\omega+\eta$，其中 $\eta$ 是闭形式，并且对任意
$\xi\in TM$ 都有 $\eta(\xi,\mathord\cdot)=0$。还有
\[
 (\Phi_*\Omega)|_{E|_{\partial_iS}}=\omega+d(H\,dt).
\]
考虑由 $\Phi\circ v(z)=(z,w(z))$ 给出的截面 $v:D^2\to E$。于是
\[
\begin{aligned}
 -\int_{D^2}v^*\Omega
 &=-\int_{D^2}(\Phi\circ v)^*(\Phi_*\Omega)
   =-\int_{D^2}(\Phi\circ v)^*(\omega+\eta)\\
 &=-\int_{D^2}w^*\omega
   -\int_{D^2}(\Phi\circ v)^*\eta.
\end{aligned}
\]
这里，$2$-形式 $(\Phi\circ u)^*\eta$ 是 $D^2$ 上的闭形式，并且它在
$\partial D^2$ 上限制为
$w^*d(H\,dt)=d(w^*H\,dt)$。由于
$H^1(D^2,\partial D^2)=0$，形式 $(\Phi\circ v)^*\eta$ 是正合的，因而
\[
 \int_{D^2}(\Phi\circ v)^*\eta
 =\int_{\partial D^2}w^*H\,dt
 =\int_{\partial D^2}H(t,w|_{\partial D^2}(t))\,dt.
\]
另一方面，截面 $u$ 的水平性蕴含
$(\partial w/\partial t)(t)=X_H(t,w(t))$ 对所有 $t\in\partial D^2$ 成立。这又蕴含，对某个 $x\in M$，有
$w|_{\partial D^2}(t)=\phi_H^t(x)$。证明完成。
\end{proof}

\subsection{曲率与弱耦合}
\stfhIndex{I301}{ruo^ou^he}{弱耦合（weak coupling）}
\label{subsec:curvature-weak-couplings}

设 $\omega_S$ 是 $S$ 上满足 $\int_S\omega_S=1$ 的面积形式，并设
$\Omega$ 是 $E$ 上与 $\Gamma$ 相伴的耦合形式。我们知道，对任意这样的
$\Omega$，
\[
 \omega_{\Gamma,\lambda}:=\Omega_\Gamma+\lambda\pi^*\omega_S
\]
对所有充分大的正常数 $\lambda$ 都是非退化的。当我们指定各边界分支上的和乐时，这样的常数 $\lambda$ 的选取将依赖于给定的和乐。

给定平凡化 $\Phi:E\to S\times(M,\omega)$ 后，可以把一条纤维等同于
$(M,\omega)$。于是，可把 $\pi:E\to S$ 上联络 $\Gamma$ 的曲率
$\operatorname{curv}(\Gamma)$ 等同于
$\Omega^2(S;C_0^\infty(M))$ 中的一个 $2$-形式。注意，在一个平凡化中，哈密顿联络在每一点 $z\in\partial_iS$ 处都确定了一个位于
$\Ham(M,\omega)$ 中的和乐。

对给定函数 $h\in C^\infty(M)$，记
\[
 \operatorname{osc}(h)=\max_{x\in M}h(x)-\min_{x\in M}h(x).
\]
回顾，在 $\Symp(M,\omega)$ 对 $C_0^\infty(M)$ 的作用下，
$\operatorname{osc}$ 是不变的。利用这一点，可按如下方式定义曲率
$\operatorname{curv}(\Gamma)\in\Omega^2(S,C_0^\infty(M))$ 的霍费尔
范数。对给定平凡化 $\Phi:E\to S\times M$，$\operatorname{curv}(\Gamma)$
诱导斜对称双线性形式
\[
 K_z^{\Gamma,\Phi}:T_zS\times T_zS\longrightarrow C_0^\infty(M).
\]
因而可定义它的霍费尔范数为
\begin{equation}
 \|K_z^{\Gamma,\Phi}\|:=
 \|K_z^{\Gamma,\Phi}(v,w)\|
 \quad\text{其中 }\omega_S(v,w)=1.
 \label{eq:curvature-hofer-norm}
\end{equation}
显然，只要 $\omega_S(v,w)=1$，右端便与 $(v,w)$ 的选取无关。

\begin{exercise}\label{ex:curvature-norm-trivialization-independent}
证明定义 (4.3.17) 与平凡化的选取无关。更具体地，考虑两个平凡化
$\Phi,\Psi$，并写
\[
 \Psi\circ\Phi^{-1}(z,x)
 =(z,\phi_{\Psi\Phi}^z(x)),
\]
其中 $\phi_{\Psi\Phi}^z\in\Ham(M)$。证明
\[
 K_z^{\Gamma,\Psi}(v,w)
 =K_z^{\Gamma,\Phi}(v,w)\circ\phi_{\Psi\Phi}^z.
\]
\end{exercise}

现在可以对 $E\to S$ 的某个（因而任意）平凡化定义
\[
 \|\operatorname{curv}(\Gamma)\|=\|K^{\Gamma,\Phi}\|,
\]
并称之为 $\Gamma$ 的\emph{曲率宽度}。依照 (Gr96)、(Po96) 与 (En00)，把联络 $\Gamma$ 的\emph{$K$-面积}定义为
\stfhIndex{I089}{qu^lv^kuan^du}{曲率宽度（curvature width）}\stfhIndex{I024}{00-symbol-024}{$K$-面积（$K$-area）}
\[
 K\text{-}\operatorname{area}(\Gamma)
 =\frac{1}{\|\operatorname{curv}(\Gamma)\|},
\]
并把 $E$ 的 $K$-面积定义为
\[
 K\text{-}\operatorname{area}(E)
 =\inf_\Gamma K\text{-}\operatorname{area}(\Gamma).
\]

下面的定义将在本书第三部研究三角不等式时用到。

\begin{definition}\label{def:k-area-holonomy-tuple}
把配备和乐组 $C=(C_1,\ldots,C_h)$ 的 $E$ 的 $K$-面积定义为
\begin{equation}
 K\text{-}\operatorname{area}(E;C)
 =\inf_{\Gamma\in\mathcal L(C)}K\text{-}\operatorname{area}(\Gamma).
 \label{eq:k-area-holonomy-tuple}
\end{equation}
\end{definition}

由此可知，$K\text{-}\operatorname{area}(E;C)$ 与共轭类
$C_1,\ldots,C_\ell$ 的次序无关。

现在可以陈述一个联系 $E$ 的两个不变量 $\operatorname{size}(E)$ 与
$K\text{-}\operatorname{area}(E)$ 的定理。令 $\Omega=\Omega_\Gamma$ 为
$\Gamma$ 的耦合形式。它确定了 $E$ 上一个自然的辛形式形变类，其中的形式为
\[
 \omega_{\Gamma,\lambda}:=\Omega_\Gamma+\lambda\pi^*\omega_S,
\]
这里 $\omega_S$ 是面积形式，而 $\lambda>0$ 是充分大的常数。我们总把
$\omega_S$ 归一化为 $\int_S\omega_S=1$。

下面的定义是定义 4.3.3 的直接推广。

\begin{definition}\label{def:size-with-holonomy}
考虑带有由耦合形式 $\Omega$ 诱导的联络 $\Gamma$ 的哈密顿纤维化
$E\to S$。把 $\operatorname{size}(\Gamma)$ 定义为命题 4.3.1 中出现的
$1/\lambda$ 的上确界，其中形式
$\Omega_\Gamma+\lambda\pi^*\omega_\Sigma$ 非退化。定义
\begin{equation}
 \operatorname{size}(E;C)
 =\sup_{\Gamma\in\mathcal L(C)}\operatorname{size}(\Gamma).
 \label{eq:size-with-holonomy}
\end{equation}
\end{definition}

下面的定理等价于恩托夫证明的定理 3.7.4 (En00)。

\begin{theorem}\label{thm:size-k-area}
设 $\pi:E\to\Sigma$ 是带联络 $\Gamma\in\mathcal L(C)$ 的哈密顿纤维化。则
\begin{equation}
 \operatorname{size}(E;C)
 \leq K\text{-}\operatorname{area}(E;C).
 \label{eq:size-k-area}
\end{equation}
\end{theorem}

\begin{proof}
为检验 $\Omega_\Gamma+\lambda\pi^*\omega_\Sigma$ 的非退化性，只需保证该形式在每个 $e\in E$ 的 $HT_eE$ 上非退化。但对任意给定的
$\eta_1,\eta_2\in T_{\pi(e)}S$，都有
\[
 \Omega(\eta_1^\#,\eta_2^\#)+\lambda\omega_\Sigma(\eta_1,\eta_2)
 =\operatorname{curv}(\Gamma)(\eta_1,\eta_2)
  +\lambda\omega_\Sigma(\eta_1,\eta_2).
\]
因此，若
$\operatorname{curv}(\Gamma)(\eta_1,\eta_2)
 +\lambda\omega_\Sigma(\eta_1,\eta_2)=0$ 且
$\omega_\Sigma(\eta_1,\eta_2)\ne0$，则
\[
 \lambda=-\frac{\operatorname{curv}(\Gamma)(\eta_1,\eta_2)}
 {\omega_\Sigma(\eta_1,\eta_2)}.
\]
特别地，若 $\omega_\Sigma(\eta_1,\eta_2)=1$，则
$\lambda=-\operatorname{curv}(\Gamma)(\eta_1,\eta_2)$。这蕴含，只要
\[
 \lambda>
 \max_{z\in\Sigma}
 \max_{\eta_1,\eta_2;\,\omega_\Sigma(\eta_1,\eta_2)=1}
 |\operatorname{curv}(\Gamma)(\eta_1,\eta_2)|,
\]
形式 $\omega_{\Gamma,\lambda}=\Omega_\Gamma+
\lambda\pi^*\omega_\Sigma$ 就非退化。

等价地，只要能找到哈密顿联络 $\Gamma$，使得
\[
 \frac1\lambda
 <\frac1{\displaystyle\max_{z\in\Sigma}
   \|\operatorname{curv}(\Gamma)(z)\|}
 =\frac1{\|\operatorname{curv}(\Gamma)\|},
\]
形式 $\Omega_\Gamma+\lambda\pi^*\omega_\Sigma$ 就非退化。对所有这样的
$\Gamma$，取 $1/\|\operatorname{curv}(\Gamma)\|$ 的下确界，便得到：对任意给定的 $0<\epsilon<\lambda$，只要
\[
 \frac1\lambda\leq K\text{-}\operatorname{area}(E;C)+\epsilon,
\]
形式 $\omega_{\Gamma,\lambda}$ 就非退化。对所有这样的 $\lambda$ 取
$1/\lambda$ 的上确界，再令 $\epsilon\to0$，便证明了
\[
 \operatorname{size}(E;C)
 \leq K\text{-}\operatorname{area}(E;C).
\]
证明完成。
\end{proof}

还可以让联络 $\Gamma$ 在固定一个和乐同伦类的联络之间变化，从而进一步细化对纤维化 $\pi:E\to\Sigma$ 的曲率与大小的研究。

\subsection{带标记哈密顿纤维化}
\stfhSubIndex{I158}{ha^mi^dun^xian^wei^hua}{哈密顿纤维化（Hamiltonian fibration）}{dai^biao^ji^ha^mi^dun^xian^wei^hua}{带标记哈密顿纤维化（marked Hamiltonian fibration）}
\label{subsec:marked-hamiltonian-fibration}

现在考虑带耦合形式 $\Omega$ 的哈密顿纤维化 $E\to S$，并以
$\nabla=\nabla_\Omega$ 表示相伴的辛联络。按定义，$\nabla$ 沿一条以点
$z$ 为基点的环路的和乐定义了 $\Ham(E_z,\omega)$ 中的一个元素。特别地，当
$\partial S$ 的每个边界分支都带有一个标记点 $*$ 以及一个等同
\[
 (E_*,\Omega_*)\longrightarrow(M,\omega)
\]
时，和乐就定义了 $\Ham(M,\omega)$ 中的一个元素。

\begin{definition}\label{def:marked-hamiltonian-fibration}
设 $S$ 是紧致有边黎曼曲面。记
\[
 \partial S=\bigcup_{i=1}^h=\partial_iS.
\]
\begin{enumerate}[label=(\arabic*)]
\item
带标记哈密顿纤维化 $E\to S$ 是一个哈密顿纤维化 $(E,\Omega)$，其中在每个 $\partial_iS$ 上有一个标记点 $z_i$，并有一个等同
$(E_{z_i},\Omega_{z_i})\cong(M,\omega)$。
\item
若带标记哈密顿纤维化 $E$ 满足
\begin{equation}
 \operatorname{hol}_\nabla(\partial_iS;z_i)=\phi_i,
 \label{eq:marked-boundary-holonomy}
\end{equation}
则称它具有边界和乐数据
$(\phi_1,\ldots,\phi_h)\in\Ham(M,\omega)^h$。
\end{enumerate}
\end{definition}

下面的定义对于以后用弗洛尔理论研究哈密顿几何也很有用。

\begin{definition}\label{def:marked-hamiltonian-fibration-connection}
设 $S$ 如上。一个\emph{带联络的带标记哈密顿纤维化}，是一个带有沿边界
$\partial S$ 的平凡化
\[
 \Phi_{\partial S}:
 (E|_{\partial S},\Omega|_{\partial S})
 \longrightarrow\partial S\times(M,\omega)
\]
的哈密顿纤维化 $(E,\Omega)$。若
\begin{equation}
 \operatorname{hol}_{\nabla;0}^t(\partial_iS)=\phi_{H_i}^t,
 \qquad i=1,\ldots,h,
 \label{eq:marked-connection-boundary-holonomy}
\end{equation}
则称带联络的带标记哈密顿纤维化具有边界和乐数据
$C=(C_{H_1},\ldots,C_{H_h})$。
\end{definition}

注意，若选取另一个平凡化
\[
 \Phi'_{\partial S}:
 (E|_{\partial S},\Omega|_{\partial S})
 \longrightarrow\partial S\times(M,\omega),
\]
则相应的和乐彼此共轭，因而这些共轭类不依赖于平凡化，而依赖于
$\partial S$ 中各标记点处的等同。
